/media/dominic/Data1/Box/Letters/letter_template.tex
\signature{\Who\Title\\ \makebox[0pt][l]{\small on behalf of the co-authors: Jocelyn Lavallee, Anna Raffeld, and Doria Gordon}}
\begin{document}
\begin{letter}{
		The Editors, \emph{Energy \& Environmental Science}\\
	\vspace{0.01\textheight} 
    \textbf{Re.  Submission of manuscript.}
	\vspace{-0.04\textheight} 
}

\opening{Dear editors,}

I am pleased to submit for your consideration a manuscript entitled ``Coordinating bioenergy, biochar and geological storage for biomass-based carbon dioxide removal'', for publication in \emph{Energy \& Environmental Science} as a Full Paper.

The window for limiting global warming to below 2~°C is closing fast. The world will need billions of tonnes of carbon dioxide removal (CDR) per year by mid-century, and IPCC mitigation pathways look to biomass for most of it. But, land is already stretched between food and nature, and each tonne of biomass can be used only once: combusted for electricity (bioenergy, BE), combusted with carbon capture and storage (BECCS), pyrolyzed to biochar for soil (pyrogenic carbon capture and storage, PyCCS), or the counterfactual of simply leaving it in the field to provide soil organic matter (SOM). Decisions on this resource are being made now, yet there is a critical gap in the evidence needed to inform them.

Integrated assessment models (IAMs) largely assume that engineered CDR is from BECCS and ignore biochar. They also lack the resolution to accurately quantify the spatially heterogeneous trade-offs that decide between options. Bottom-up studies capture this local detail, but evaluate individual technologies without allocating resources efficiently between them. Policy makers are therefore planning large-scale removal without knowing where, when, or at what carbon price each technology is the better use of the same biomass.

Our manuscript fills this gap. We analyze the GHG abatement, costs, and co-benefits of using crop residues (an abundant biomass resource that needs no additional land) for BE, BECCS, and PyCCS at 0.1° resolution across the United States, Europe, China, and India, four of the largest agricultural economies, with sharply contrasting soils, energy markets, and geological storage. CO$_{2}$ is routed along least-cost pipelines to characterized storage basins, accounting for terrain, protected areas, multimodal transport and injection costs. Biochar is valued through its effects on carbon sequestration, crop production, and soil GHG emissions. 

We find that PyCCS is the cheapest removal at moderate carbon prices, and in parts of China and India is profitable with no carbon price at all. BECCS stores more than twice as much carbon per tonne of biomass and takes over at higher carbon prices, tripling abatement to about 2 Gt CO$_{2}$e per year. The crossover price from PyCCS to BECCS varies with soils, crops, energy prices, biomass supply logistics, and routed distance to storage, so no single supply curve can say which technology belongs where. Although cost-effective at lower carbon prices, building PyCCS first as carbon prices rise risks locking in residue that BECCS could use better where shared CO$_{2}$ infrastructure exists; where it does not, PyCCS is the low-regret first step.

These findings both challenge the BECCS-centered view of CDR, and also highlight the importance of detailed spatially explicit models to inform local practice. In its optimal locations, biochar can start removing carbon now, while rewarding farmers with better harvests; BECCS can then deliver larger removals once CO$_{2}$ infrastructure is in place. Bringing together energy economics, process engineering, biogeochemistry, and geospatial analysis, we believe this work will be of great interest to readers of \emph{Energy \& Environmental Science} and to the wider scientific and policy community designing and analyzing climate change mitigation pathways.

The manuscript is original, has not been published previously, and is not under consideration elsewhere. All authors have approved the submission, and there are no conflicts of interest. We suggest the following people as potential reviewers: \textbf{Alexander Popp}, Potsdam Institute for Climate Impact Research (PIK), popp@pik-potsdam.de (expertise in integrated assessment modeling, land and biomass for climate change mitigation). \textbf{Annette Cowie}, NSW Department of Primary Industries, annette.cowie@dpird.nsw.gov.au (expertise in bioenergy, biochar, life cycle assessment, and GHG accounting). \textbf{Vassilis Daioglou}, PBL Netherlands Environmental Assessment Agency and Utrecht University, vassilis.daioglou@pbl.nl (expertise in integrated assessment modeling of biomass supply, bioenergy, and BECCS). \textbf{Pete Smith}, University of Aberdeen, pete.smith@abdn.ac.uk (expertise in soil carbon, biochar, BECCS, and land-based climate change mitigation).

We thank you for your consideration and look forward to your response.

\closing{Yours sincerely,}

\end{letter}
\end{document}
